% Options for packages loaded elsewhere
\PassOptionsToPackage{unicode}{hyperref}
\PassOptionsToPackage{hyphens}{url}
\documentclass[
  12pt,
]{article}
\usepackage{xcolor}
\usepackage{amsmath,amssymb}
\setcounter{secnumdepth}{-\maxdimen} % remove section numbering
\usepackage{iftex}
\ifPDFTeX
  \usepackage[T1]{fontenc}
  \usepackage[utf8]{inputenc}
  \usepackage{textcomp} % provide euro and other symbols
\else % if luatex or xetex
  \usepackage{unicode-math} % this also loads fontspec
  \defaultfontfeatures{Scale=MatchLowercase}
  \defaultfontfeatures[\rmfamily]{Ligatures=TeX,Scale=1}
\fi
\usepackage{lmodern}
\ifPDFTeX\else
  % xetex/luatex font selection
\fi
% Use upquote if available, for straight quotes in verbatim environments
\IfFileExists{upquote.sty}{\usepackage{upquote}}{}
\IfFileExists{microtype.sty}{% use microtype if available
  \usepackage[]{microtype}
  \UseMicrotypeSet[protrusion]{basicmath} % disable protrusion for tt fonts
}{}
\makeatletter
\@ifundefined{KOMAClassName}{% if non-KOMA class
  \IfFileExists{parskip.sty}{%
    \usepackage{parskip}
  }{% else
    \setlength{\parindent}{0pt}
    \setlength{\parskip}{6pt plus 2pt minus 1pt}}
}{% if KOMA class
  \KOMAoptions{parskip=half}}
\makeatother
\usepackage{longtable,booktabs,array}
\usepackage{calc} % for calculating minipage widths
% Correct order of tables after \paragraph or \subparagraph
\usepackage{etoolbox}
\makeatletter
\patchcmd\longtable{\par}{\if@noskipsec\mbox{}\fi\par}{}{}
\makeatother
% Allow footnotes in longtable head/foot
\IfFileExists{footnotehyper.sty}{\usepackage{footnotehyper}}{\usepackage{footnote}}
\makesavenoteenv{longtable}
\setlength{\emergencystretch}{3em} % prevent overfull lines
\providecommand{\tightlist}{%
  \setlength{\itemsep}{0pt}\setlength{\parskip}{0pt}}
\usepackage{bookmark}
\IfFileExists{xurl.sty}{\usepackage{xurl}}{} % add URL line breaks if available
\urlstyle{same}
\hypersetup{
  hidelinks,
  pdfcreator={LaTeX via pandoc}}

\author{SCX}
\date{2026}

\begin{document}

\section{Cercis Algorithm}\label{cercis-algorithm}

\begin{quote}
\textbf{Formal name}: Cercis Self-Evolving Gatekeeper (CESG)
\textbf{Location}:
/g/Xiaogan\_Supercomputing\_data/SCX/theory/self\_evolution/
\textbf{Paper target}: Nature Computational Science (Paper 2 of
two-paper strategy)
\end{quote}

\begin{center}\rule{0.5\linewidth}{0.5pt}\end{center}

\subsection{Why ``Cercis''}\label{why-cercis}

\emph{Cercis chinensis} is a \textbf{cauliflorous} tree --- its
flowers bloom directly from old branches and trunks, not from new
shoots.

\begin{verbatim}
Cercis biological metaphor → SCX mathematical reality:

  Old branches & trunk    →  Memory bank M_t (accumulated structures, never deleted)
  Flowers                 →  Gatekeeper evaluations (new scores on old structures)
  Each spring bloom       →  Each iteration: S_t re-scores M_t
  Cauliflory              →  Resurrection: discarded structures "bloom" when S_t matures
  New shoots              →  Novelty bonus η(t) · N_t(s) — exploration
  Deep roots              →  DFT/experimental ground truth (physical anchor)
\end{verbatim}

\textbf{Tagline for paper}: \emph{``Knowledge does not grow from new
data alone --- it flowers from the old wood of accumulated
experience.''}

\begin{center}\rule{0.5\linewidth}{0.5pt}\end{center}

\subsection{Paper 2: Cercis Flowers}\label{paper-2-cercis-flowers}

\begin{longtable}[]{@{}
  >{\raggedright\arraybackslash}p{(\linewidth - 2\tabcolsep) * \real{0.4286}}
  >{\raggedright\arraybackslash}p{(\linewidth - 2\tabcolsep) * \real{0.5714}}@{}}
\toprule\noalign{}
\begin{minipage}[b]{\linewidth}\raggedright
Item
\end{minipage} & \begin{minipage}[b]{\linewidth}\raggedright
Detail
\end{minipage} \\
\midrule\noalign{}
\endhead
\bottomrule\noalign{}
\endlastfoot
\textbf{Title} & \emph{Cercis Flowers: A Self-Evolving Gatekeeper with
Provable Convergence} \\
\textbf{Core Theorems} & Theorem SE-1 (Convergence), Theorem SE-2
(Completeness Bound) \\
\textbf{Target} & Nature Computational Science / Nature Machine
Intelligence \\
\textbf{Depends on} & SCX framework (Paper 1, JMLR/TMLR) \\
\textbf{Location} & \texttt{theory/self\_evolution/} (theorems),
\texttt{paper/paper\_gatekeeper/} (manuscript) \\
\end{longtable}

\end{document}
